\begin{pastexam}[label={px:df3-096-ds-02}]{地方三等資訊處理 96 年 第二題}
（申論；20 分）

令一疊紙牌依序為 $a_1, a_2, \ldots, a_n$。洗牌動作 shuffle($q$), $q < n$,為將紙牌分成
$a_1, a_2, \ldots, a_q$ 與 $a_{q+1}, a_{q+2}, \ldots, a_n$ 兩疊，再將兩疊紙牌逐張交錯放置成一疊為
$a_1, a_{q+1}, a_2, a_{q+2}, \ldots, a_q, a_{2q}, a_{2q+1}, a_{2q+2}, \ldots, a_n$ 當 $q \le n/2$,
$a_1, a_{q+1}, a_2, a_{q+2}, \ldots, a_{n-q}, a_n, a_{n-q+1}, a_{n-q+2}, \ldots, a_q$ 當 $q > n/2$。
用單向鏈結串列（single-linked list）來模擬此疊紙牌。試寫一段程式模擬洗牌動作
shuffle($q$)。（20 分）

範例：一疊紙牌為
\begin{center}
\begin{tikzpicture}[x=1.45cm, >=Stealth,
  d/.style={draw, minimum width=8.5mm, minimum height=6.5mm, inner sep=1pt},
  p/.style={draw, minimum width=2.6mm, minimum height=6.5mm, inner sep=0pt, anchor=west}]
\node (h) at (0,0) {head};
\foreach \c [count=\i] in {$\clubsuit$5, $\diamondsuit$K, $\spadesuit$9, $\clubsuit$7, $\diamondsuit$2, $\heartsuit$8} {
  \node[d] (d\i) at (\i,0) {\c};
  \node[p] (p\i) at (d\i.east) {};
}
\node (nil) at (7.15,0) {Nil};
\draw[->] (h) -- (d1);
\foreach \i [evaluate=\i as \j using int(\i+1)] in {1,...,5} \draw[->] (p\i.center) -- (d\j.west);
\draw[->] (p6.center) -- (nil);
\end{tikzpicture}
\end{center}
則經洗牌動作 shuffle(3)後為
\begin{center}
\begin{tikzpicture}[x=1.45cm, >=Stealth,
  d/.style={draw, minimum width=8.5mm, minimum height=6.5mm, inner sep=1pt},
  p/.style={draw, minimum width=2.6mm, minimum height=6.5mm, inner sep=0pt, anchor=west}]
\node (h) at (0,0) {head};
\foreach \c [count=\i] in {$\clubsuit$5, $\clubsuit$7, $\diamondsuit$K, $\diamondsuit$2, $\spadesuit$9, $\heartsuit$8} {
  \node[d] (d\i) at (\i,0) {\c};
  \node[p] (p\i) at (d\i.east) {};
}
\node (nil) at (7.15,0) {Nil};
\draw[->] (h) -- (d1);
\foreach \i [evaluate=\i as \j using int(\i+1)] in {1,...,5} \draw[->] (p\i.center) -- (d\j.west);
\draw[->] (p6.center) -- (nil);
\end{tikzpicture}
\end{center}
\end{pastexam}

\begin{pastsoln}
原卷的方塊、紅心印成實心符號，這裡用 TeX 的標準花色符號 $\diamondsuit$、$\heartsuit$ 排版，不影響題意。
題目只說 $q < n$，沒有說 $q \ge 1$；$q \le 0$ 時不用動，程式裡一併處理。

\textbf{想法。} 兩種情況其實是同一件事：兩疊輪流各取一張，\textbf{哪一疊先用完，另一疊剩下
的就整串接在後面}。 $q \le n/2$ 時第一疊先用完，剩下 $a_{2q+1}, \ldots, a_n$；$q > n/2$ 時
第二疊先用完，剩下 $a_{n-q+1}, \ldots, a_q$。 用串列做不必搬資料，只要改指標：
\begin{enumerate}[nosep]
\item 從 head 走 $q - 1$ 步到 $a_q$，$a_{q+1}$ 就是第二疊的開頭；把 $a_q$ 的 next 設成 NULL，
切成兩條串列。
\item $p$ 指第一疊目前這張、$s$ 指第二疊目前這張。 先記下兩者原本的下一張 $p'$、$s'$，再改成
$p \to s \to p'$，然後 $p \gets p'$、$s \gets s'$。
\item $p'$ 是 NULL（第一疊用完）時，$s$ 後面原本就串著第二疊剩下的牌，不用再動，直接結束；
$s'$ 是 NULL（第二疊用完）時，$s \to p'$ 已經把第一疊剩下的牌接上，迴圈也自然結束。
\end{enumerate}

\textbf{程式（C）。}
\begin{minted}[fontsize=\small, style=perldoc]{c}
typedef struct Node {
    char card[4];            /* 例如 "C5"、"DK" */
    struct Node *next;
} Node;

/* 對 head 開頭的串列做 shuffle(q)，第一張仍是 a1，所以 head 不變 */
void shuffle(Node *head, int q)
{
    if (head == NULL || q <= 0) return;
    Node *p = head;
    for (int k = 1; k < q && p != NULL; k++)   /* p 走到第 q 張 a_q */
        p = p->next;
    if (p == NULL || p->next == NULL) return;  /* q >= n：沒有第二疊 */
    Node *s = p->next;          /* 第二疊的第一張 a_{q+1} */
    p->next = NULL;             /* 切成兩疊 */
    p = head;
    while (p != NULL && s != NULL) {
        Node *pn = p->next;     /* 第一疊的下一張 */
        Node *sn = s->next;     /* 第二疊的下一張 */
        p->next = s;            /* a_i 後面接 a_{q+i} */
        if (pn == NULL) break;  /* 第一疊用完，第二疊剩下的已經接在 s 後面 */
        s->next = pn;           /* a_{q+i} 後面接 a_{i+1} */
        p = pn;
        s = sn;
    }
}
\end{minted}

\textbf{追蹤範例 shuffle(3)}，$n = 6$。 $p$ 走到第 $3$ 張 $\spadesuit 9$，第二疊從
$\clubsuit 7$ 開始；切開後兩疊是 $\clubsuit 5, \diamondsuit K, \spadesuit 9$ 和
$\clubsuit 7, \diamondsuit 2, \heartsuit 8$。
\begin{center}
\begin{tabular}{cllll}
\toprule
輪 & $p$ & $s$ & 改的指標 & 目前從 head 串起來的前段 \\
\midrule
1 & $\clubsuit 5$ & $\clubsuit 7$ & $\clubsuit 5 \to \clubsuit 7 \to \diamondsuit K$ & $\clubsuit 5, \clubsuit 7, \diamondsuit K$ \\
2 & $\diamondsuit K$ & $\diamondsuit 2$ & $\diamondsuit K \to \diamondsuit 2 \to \spadesuit 9$ & $\ldots, \diamondsuit K, \diamondsuit 2, \spadesuit 9$ \\
3 & $\spadesuit 9$ & $\heartsuit 8$ & $\spadesuit 9 \to \heartsuit 8$；$p'$ 是 NULL，結束 & $\ldots, \spadesuit 9, \heartsuit 8$ \\
\bottomrule
\end{tabular}
\end{center}
結果 $\clubsuit 5, \clubsuit 7, \diamondsuit K, \diamondsuit 2, \spadesuit 9, \heartsuit 8$，
和題目的範例相同。

\textbf{另外兩種情況}：$n = 6$、$q = 2$（$q < n/2$）：兩疊 $a_1 a_2$、$a_3 a_4 a_5 a_6$，第二輪
$p' = $ NULL 時 $a_4$ 後面仍串著 $a_5, a_6$，得 $a_1, a_3, a_2, a_4, a_5, a_6$。 $q = 4$
（$q > n/2$）：兩疊 $a_1 \ldots a_4$、$a_5 a_6$，第二輪做完 $a_6 \to a_3$ 後 $s' = $ NULL，迴圈
結束，$a_3 \to a_4 \to$ NULL 原封不動，得 $a_1, a_5, a_2, a_6, a_3, a_4$。 兩者都符合題目的
公式。

\textbf{複雜度}：找 $a_q$ 走 $q - 1$ 步，交錯每輪 $O(1)$、最多 $\min(q, n - q)$ 輪，合計
$O(n)$ 時間；只用了幾個指標，額外空間 $O(1)$。 串列的基本操作見資結書第~5.1~節。
\end{pastsoln}
